We have presented BibleGoAI, a multi-stage pipeline for synthesizing concise, source-adherent summaries of biblical commentary using large language models, retrieval-augmented generation, and semantic similarity scoring. The pipeline was run once over a curated corpus of 14 commentators spanning four centuries, producing Brief Helpful Texts (BHTs) of roughly 80 words each (see Section~\ref{sec:evaluation} for per-pipeline distributions) that are grounded in source material, quality-scored across multiple dimensions, and traceable via automated footnotes. The resulting BHTs and their source commentaries are served through BibleGo.org, reaching approximately 100--300 unique users per month at the time of writing.

\paragraph{Future work.}
Several directions remain for future development:
\begin{itemize}[nosep,leftmargin=*]
    \item \textbf{Regeneration with stronger models}: Re-running the pipeline against the same source corpus using more recent frontier LLMs to measure how much of the residual quality gap is attributable to the GPT-3.5-era generator versus the pipeline design itself.
    \item \textbf{Live admin-driven regeneration}: Deploying the pipeline as a live service so that administrators can regenerate individual BHTs (or whole books) on demand as commentators are added, prompts are refined, or source-text corrections land, rather than treating BHT production as a one-shot offline process.
    \item \textbf{Publishing the digitized commentaries}: Releasing the digitized texts of Henry Alford's \textit{New Testament for English Readers} and Robert Govett's \textit{Commentary on Revelation} as standalone open resources, independent of the BHT pipeline, so the human digitization effort is reusable beyond this project.
    \item \textbf{Additional commentators}: Continuing to expand the curated corpus with further historical and contemporary commentators; Govett's commentary on Revelation was completed and added to the digitized corpus in 2025 following the work reported here, and similar additions are under consideration.
    \item \textbf{Enhanced original-language tools}: Considered for future iterations, expanding word-study features beyond Strong's concordance to include Hebrew morphology, cross-references, and parallel passage analysis.
\end{itemize}

\paragraph{Lessons for similar projects.}
Four patterns from this work generalize beyond biblical commentary. First, an extractive intermediate representation---the \emph{choicest quotes} phase---proved decisive: it gives every downstream stage an auditable artifact and makes source fidelity checkable without re-reading every commentary. Second, a best-of-$N$ loop against a hard constraint gate is a low-cost way to get acceptable coverage from an imperfect generator, at the price of a small fraction of retained outputs that miss the nominal quality threshold; the attempt-level analysis in Section~\ref{sec:evaluation} shows this trade-off concretely. Third, a volunteer-driven curation and digitization effort can substantially expand a domain's source corpus beyond what any individual team could assemble, but only when paired with a structured consolidation pipeline that merges human correction back into machine-readable form; the same digitization pipeline served both the Alford and Govett efforts, demonstrating its replicability across projects. Fourth, automated footnote generation against the source corpus proved decisive for credibility and verifiability. Beyond the engineering value of similarity scoring, the visible per-claim footnote chain in the served BHTs gave readers a tangible verification path: the difference between ``trust the LLM'' and ``click to read where this came from.'' For domains where source-claim grounding determines whether users will trust the system at all (medical, legal, theological), this is not optional polish; it is the feature that makes the system usable. Domains with similar characteristics (specialized corpora, community stewardship, and a need for claim-level traceability) may find the overall pattern reusable.
